\documentclass{article}
\usepackage{amsmath}
\usepackage[utf8]{inputenc}
\usepackage[brazil]{babel}
\usepackage{geometry}
\geometry{a4paper, margin=2cm}

\begin{document}

\section*{Resolução da Questão 107 --- ENEM 2024 --- Ciências da Natureza}

\textbf{Alternativa correta: E}

A alternativa E está correta porque corresponde ao método que realmente remove o mercúrio contaminante do solo, acumulando-o na planta e permitindo a retirada da biomassa contaminada, impedindo assim a entrada do mercúrio na cadeia alimentar.

A estratégia adotada para resolver essa questão inclui a leitura atenta do enunciado e a análise da figura explicativa dos métodos de fitorremediação. O foco está em entender qual técnica remove o mercúrio do ambiente de forma definitiva, comparando as descrições dos processos.

\begin{enumerate}
  \item Identificar o objetivo: remoção do mercúrio evitando sua entrada na cadeia alimentar.
  \item Analisar a figura e as definições: fitoestabilização, fitovolatilização, fitodegradação, fitoestimulação e fitoextração.
  \item Reconhecer que a fitoextração envolve o acúmulo do contaminante na parte aérea das plantas e posterior retirada da biomassa.
  \item Compreender que as outras técnicas não promovem a remoção do mercúrio, mas sua imobilização, transformação ou estímulo microbiano que não retiram o metal do solo.
  \item Confirmar que a alternativa E corresponde exatamente ao objetivo pedido.
\end{enumerate}

\section*{Explicação Didática}

\subsection*{Resumo}
A fitorremediação engloba técnicas que utilizam plantas para tratar solo contaminado. Para remover o mercúrio e evitar sua entrada na cadeia alimentar, a técnica correta é a fitoextração, que acumula o metal nas partes aéreas das plantas que depois são removidas do local, eliminando o contaminante.

\subsection*{Passo a passo para compreensão}
\begin{enumerate}
  \item Entenda que o mercúrio é um metal pesado perigoso e que sua remoção física evita contaminação ambiental.
  \item Analise as definições dos métodos na figura:
  \begin{itemize}
    \item \textbf{Fitoestabilização:} imobiliza o contaminante no solo, sem removê-lo.
    \item \textbf{Fitovolatilização:} transforma o contaminante em forma volátil e o liberta no ar.
    \item \textbf{Fitodegradação:} degrada contaminantes orgânicos via enzimas, não metais pesados.
    \item \textbf{Fitoestimulação:} estimula microrganismos para degradar contaminantes orgânicos.
    \item \textbf{Fitoextração:} absorve e acumula o contaminante na biomassa aérea, permitindo sua retirada.
  \end{itemize}
  \item Compare e elimine as opções que não promovem a remoção efetiva do mercúrio.
  \item Conclua que a fitoextração (E) é a solução correta para remoção do mercúrio contaminante.
\end{enumerate}

\subsection*{Armadilhas comuns}
\begin{itemize}
  \item Confundir imobilização com remoção.
  \item Supor que a fitovolatilização elimina o contaminante definitivamente.
  \item Pensar que metais pesados podem ser degradados enzimaticamente.
  \item Acreditar que a fitoestimulação remove diretamente metais pesados.
  \item Ignorar a necessidade de retirar fisicamente a biomassa contaminada para impedir contaminação.
\end{itemize}

\section*{Revisão dos Conceitos}

\begin{tabular}{|l|p{6cm}|p{6cm}|}
\hline
\textbf{Conceito} & \textbf{Descrição} & \textbf{Aplicação no contexto} \\
\hline
Fitoextração & Remoção do poluente pela absorção e acúmulo nas partes aéreas da planta & Permite a retirada física do mercúrio, evitando sua entrada na cadeia alimentar \\
\hline
Fitoestabilização & Imobilização do contaminante no solo & Reduz disponibilidade, mas não remove o mercúrio \\
\hline
Fitovolatilização & Conversão do contaminante em forma volátil para liberação no ar & Pode causar contaminação atmosférica, não seguro para mercúrio \\
\hline
Fitodegradação & Degradação de contaminantes orgânicos via enzimas & Não aplicável a metais pesados como mercúrio \\
\hline
Fitoestimulação & Estímulo de microrganismos para degradar contaminantes orgânicos & Não remove metais pesados diretamente \\
\hline
\end{tabular}

\section*{Resolução e "Pulo do Gato"}

O problema pede o método que retira o mercúrio do solo evitando sua entrada na cadeia alimentar. A análise dos métodos mostra:

\begin{itemize}
  \item Fitoextração é a única que remove fisicamente o mercúrio acumulando-o na planta, permitindo a retirada da biomassa contaminada.
  \item Os demais métodos apenas imobilizam, transformam ou estimulam a degradação de orgânicos, não metais pesados.
\end{itemize}

\textbf{Pulo do gato:} metais pesados, como mercúrio, não são degradados em compostos menos tóxicos via enzimas; portanto, a remoção física (fitoextração) é essencial para evitar contaminação ambiental e entrada na cadeia alimentar.

\section*{Armadilhas e Distratores}

\begin{itemize}
  \item\textbf{A - Fitoestabilização}: confunde-se imobilização com remoção. Embora o mercúrio fique preso no solo, pode ainda contaminar a cadeia alimentar.
  \item\textbf{B - Fitovolatilização}: pode parecer eficaz pois remove o mercúrio do solo, mas o libera no ar, podendo causar poluição atmosférica e contaminação.
  \item\textbf{C - Fitodegradação}: incorretamente considera que mercúrio pode ser degradado como contaminantes orgânicos, o que é falso.
  \item\textbf{D - Fitoestimulação}: acredita que microrganismos estimulados removem metais pesados diretamente, o que se aplica apenas a organismos orgânicos.
\end{itemize}

\section*{Código da Questão}

\texttt{CN\_FITO\_001}

\bigskip

\textbf{Referência oficial:}
\begin{quote}
LUIZ, E. P. \textit{Avaliação dos métodos de fitorremediação [...] na remoção de chumbo, cobre e zinco.} São Paulo: UFABC, 2016 (adaptado).
\end{quote}

\end{document}